\chapter{METHODOLOGY AND IMPLEMENTATION WORKFLOW}

Proper implementation of adsorption studies is important for obtaining reliable fluoride removal data and adsorption analysis. The project involves preparation of \gls{gtmac}-modified coconut fiber, execution of batch adsorption experiments, and \gls{doe}-based adsorption studies under controlled laboratory conditions. Different adsorption parameters such as pH, adsorbent dose, fluoride concentration, and contact time are varied according to the \gls{ccd} experimental matrix. This chapter presents the detailed methodology and implementation workflow followed for fluoride adsorption studies using \gls{gtmac}-functionalized coconut fiber. 
\section{Methodology}
The methodology adopted for the present work consists of the following steps:
\begin{enumerate}
\item Collection of Raw Material: Raw coconut fiber is collected and used as the base adsorbent material for fluoride adsorption studies.  
\item Washing and Drying of Biomass: The collected coconut fiber is washed thoroughly using distilled water to remove dust particles and impurities. The material is then dried in a hot air oven at approximately 80 °C to remove moisture content.
\item Grinding and Sieving: The dried coconut fiber is ground into fine powder using a mechanical grinder. Uniform particle size is obtained using a 125 BSS sieve for adsorption studies.    
\item NaOH Pretreatment: The powdered biomass is treated using 10 percent NaOH solution for alkaline activation. The pretreatment process removes lignin and impurities while exposing cellulose hydroxyl groups required for surface modification. The treated biomass is washed repeatedly until neutral pH is obtained.    
\item GTMAC Functionalization: Surface modification is carried out using Glycidyl Trimethyl Ammonium Chloride (GTMAC) under alkaline conditions. The reaction is maintained at approximately 60 °C for 5 hours with sonication. GTMAC introduces positively charged quaternary ammonium groups onto the adsorbent surface, improving fluoride adsorption through electrostatic attraction.    
\item Drying and Storage of Modified Adsorbent: After functionalization, the modified adsorbent is washed thoroughly using distilled water and dried in a hot air oven. The prepared adsorbent is stored in airtight containers for adsorption studies.    
\item Characterization Studies: Fourier Transform Infrared Spectroscopy (FTIR) is used to identify functional groups present on the adsorbent surface and confirm GTMAC modification. Scanning Electron Microscopy (SEM) is used to study the surface morphology and structural changes after chemical modification.    
\item Preparation of Fluoride Solution: Fluoride stock solution is prepared using distilled water. Different fluoride concentrations are prepared according to the CCD experimental design.  
\item Batch Adsorption Studies: A fixed quantity of GTMAC-modified 
coconut fiber is added to fluoride solution under controlled experimental 
conditions. The adsorption studies are conducted by varying:
\begin{itemize}
	\item pH
	\item Adsorbent dose
	\item Initial fluoride concentration
	\item Contact time
\end{itemize}
\item Agitation and Adsorption Process: The adsorption mixtures are agitated using an orbital shaker for a specified contact time to ensure proper interaction between fluoride ions and the adsorbent surface.   
\item Filtration and Fluoride Measurement: After adsorption, the solution 
is filtered and residual fluoride concentration is measured using a Fluoride Ion-Selective Electrode (F$^{-}$ ISE).  
\item CCD Experimental Runs: Approximately 30 experimental runs are conducted according to the Central Composite Design (CCD) matrix generated using Design-Expert software.
\item \textbf{Adsorption Calculations:} The percentage fluoride removal is calculated using:

\begin{equation}
	\% \ Removal = \frac{C_0 - C_e}{C_0} \times 100
\end{equation}

The adsorption capacity is calculated using:

\begin{equation}
	q_e = \frac{(C_0 - C_e)V}{m}
\end{equation}

where:

\begin{itemize}
	\item $C_0$ = Initial fluoride concentration (mg/L)
	\item $C_e$ = Equilibrium fluoride concentration (mg/L)
	\item $V$ = Volume of solution (L)
	\item $m$ = Mass of adsorbent used (g)
	\item $q_e$ = Adsorption capacity (mg/g)
\end{itemize}

\item \textbf{Data Analysis and Optimization:} The adsorption data obtained from all experimental runs are tabulated and analyzed using Design-Expert software for adsorption response analysis and optimization studies.      
\end{enumerate}

This chapter presented the detailed methodology and implementation workflow followed for fluoride adsorption studies using GTMAC-modified coconut fiber. The preparation of adsorbent, GTMAC functionalization process, characterization studies, batch adsorption experiments, CCD experimental runs, and adsorption analysis procedure were discussed in detail. The adsorption results and optimization analysis obtained from these studies are presented in the next chapter. 
%\vspace{0.75cm}
%
% \textbf{The chapters should not end with figures, instead bring the paragraph explaining about the figure at the end followed by a summary paragraph.}
%
%After elaborating the various sections of the chapter (From Chapter 2 onwards), a summary paragraph should be written discussing the highlights of that particular chapter. This summary paragraph should not be numbered separately. This paragraph should connect the present chapter to the next chapter. 

